\section{Introduction}
A graph neural network combines node features with information from connected nodes to make predictions. An ensemble combines several such predictors, allowing one member to compensate for another's mistakes. The purpose of a learned ensemble is to improve its combined prediction. Its members can supply correct decisions on different examples, but averaging can also discard a correct alternative. Predictive quality therefore depends on what the members learn and how their outputs combine. We study graph neural network ensembles whose members share learned propagation weights, and compare their predictions with capable independent models.

GNNM retains four prediction paths through a common graph stack. Small member-specific factors change the input and output maps, so the shared graph layers receive different hidden states and can produce different predictions. This boundary factorization uses established BatchEnsemble components \citep{wen2020batchensemble}. Related factorized models have been studied for tabular prediction \citep{gorishniy2025tabm}, and BatchEnsemble has also been applied inside GNN layers \citep{kim2023efficient}. Our contribution is an empirical study of the predictions learned under these sharing constraints. Sharing reduces stored weights, while each member still performs graph propagation; predictive quality and execution cost are measured separately.

The main benchmark covers Roman Empire, Amazon Ratings, Minesweeper, Questions and Tolokers \citep{platonov2023critical}, with eight backbones and ten official splits per dataset. Its recorded means favor GNNM over four independent ordinary models in 28 of 40 backbone--dataset comparisons. SAGE and both graph-transformer variants favor GNNM on every dataset, while other backbones include losses. The benchmark compares each method's intended training and validation-selection procedure. Separate matched studies examine sharing within a specified recipe.

Those studies show different paths to a better combined prediction. In the matched Roman Empire SAGE comparison, shared propagation improves individual members despite lower disagreement. A later internal-factor SAGE study on encountered WikiCS development instead obtains higher pooled accuracy with weaker average members and broader coverage of correct alternatives. Its calibrated negative log-likelihood remains higher than that of the genuine independent ensemble. We report the accuracy benefit and confidence tradeoff together, with capable single, joint-untied and independent-ensemble controls.

The evidence supports sharing as a prediction-quality option whose value depends on the graph, backbone and training procedure. It does not identify one general optimization mechanism or a universally preferred factor placement. The original five-dataset results remain the central benchmark; the later studies distinguish member quality, useful alternatives and pooling behavior under their own stated protocols.

